\section{Определения и обозначения}
\label{sec:defs}

\textbf{Нейронная сеть}~--- упорядоченная четвёрка $\langle V,E,C,D\rangle$, где $V$~--- множество вершин графа, $E$~--- множество рёбер, $C$~--- множество соответствий между нейронами и элементами графа, $D$~--- множество нейронов и их параметров.

\textbf{Слой}~--- множество нейронов, на которые сигнал приходит одновременно из одного источника.

\textbf{Синаптическая функция}~--- сумма входных сигналов, умноженных на веса, минус порог:
\begin{equation}
	S_j = \sum_i x_i w_{ij} - T_j.
\end{equation}
В рассматриваемой модели пороги отсутствуют, поэтому $S_j = \sum_i x_i w_{ij}$. Выход нейрона $y_j = F(S_j)$; в линейной рециркуляционной сети $F(S)=S$.

\textbf{Обучение с учителем}~--- режим, при котором для каждого входного образа задан эталонный выход, а веса корректируются в направлении уменьшения отклонения от эталона. В автоассоциативной постановке эталоном служит сам входной вектор.

\textbf{Блок изображения}~--- прямоугольная область $r \times m$ пикселей, вырезанная из изображения с шагом $s_h$ по вертикали и $s_w$ по горизонтали.

\textbf{Образ} $x_k$~--- вектор-строка длины $n = r\,m\,|S|$, полученный последовательной записью значений каналов одного блока. Матрица $X \in \mathbb{R}^{Q \times n}$ содержит $Q$ образов по строкам.

\textbf{Кодер} $W$ имеет размер $n \times p$ и формирует скрытое представление $y_k = x_k W$ длины $p$. \textbf{Декодер} $W'$ имеет размер $p \times n$ и восстанавливает вектор $\hat{x}_k = y_k W'$.

\textbf{Покомпонентная ошибка} $e_k = \hat{x}_k - x_k$ имеет длину $n$. \textbf{Суммарная квадратическая ошибка} по выборке равна
\begin{equation}
	E = \frac{1}{2}\sum_{k=1}^{Q}\|e_k\|^2
	  = \frac{1}{2}\sum_{k=1}^{Q}\sum_{i=1}^{n}\bigl(\hat{x}_{k,i}-x_{k,i}\bigr)^2 .
	\label{eq:loss}
\end{equation}

\textbf{Эпоха}~--- один полный проход по всем $Q$ образам. Число обновлений весов за эпоху равно $Q$, поэтому полное число обновлений $U=QT$ отличается от числа эпох $T$.

\textbf{Градиент} скалярной функции потерь по матрице весов~--- матрица того же размера, составленная из частных производных:
\begin{equation}
	\nabla_W E = \left[\frac{\partial E}{\partial W_{ij}}\right].
\end{equation}

\textbf{Стохастический градиентный спуск} вычисляет поправку по одному образу $x_k$, а не по всей выборке.

\textbf{Адаптивный коэффициент обучения}~--- шаг, зависящий от текущего образа:
\begin{equation}
	\alpha_k = \frac{\beta}{\|x_k\|^2 + \|y_k\|^2},
	\label{eq:alpha}
\end{equation}
где $\beta>0$~--- безразмерный масштаб, который задаёт пользователь. При $\beta=1$ шаг нормируется полной энергией входа и скрытого представления, что делает обновление масштабно-инвариантным.

\textbf{Коэффициент сжатия} $Z(Q,n,p)$ определяется формулой~\eqref{eq:Z}. Дополнительно в работе рассчитывается байтовый коэффициент $K$, учитывающий, что исходные каналы хранятся в \texttt{uint8}, а коды и декодер~--- в \texttt{float64}.

\textbf{MSE} нормированных блоков равна $\mathrm{MSE}_n = 2E/(Qn)$; MSE готового BMP вычисляется отдельно в шкале $0\ldots255$ после денормализации и сборки изображения.
